\section{Remove common normal multiplicity before orbit reduction}
\label{sec:normal-multiplicity}

The previous section provides independently replayable normal-topology
proofs. Its measurements also expose a remaining avoidable cost: scaling a
fixed normal vector enlarges every endpoint in the three interval queries
and every recorded endpoint in the proofs. We now divide out common
coordinate multiplicity before those queries, then recover the original
surface's topology by an exact lifting formula. This improves computation
and proof size on repeated surfaces. The source triangulation and full
input vector are still validated before reduction.

\subsection{Why scaling is not always a disjoint union}

Let $y$ be an admissible integral normal vector in the supported orientable
three-manifold, and let $x=ky$ for a positive integer $k$. The normal
surface for $x$ can be placed in the normal interval neighbourhood of the
surface for $y$. Over each normal disc it has $k$ ordered local sheets.
Across a face the sheet order is either preserved or reversed, so the
monodromy on sheet indices is generated by
\[
 j\longmapsto k-1-j,\qquad 0\le j<k.
\]
There are no arbitrary permutations of this ordered normal fibre.
In the maintained interval model the same fact follows directly by writing
an edge-point index as $kq+j$: a scaled translation preserves $j$ and a
scaled reflection replaces it by $k-1-j$.

A two-sided component has trivial normal monodromy. It therefore contributes
$k$ disjoint copies of itself. A one-sided component has nontrivial
monodromy: its sheet pairs $\{j,k-1-j\}$ each form a connected normal
orientation double cover, with a central original component only when
$k$ is odd. Since the ambient manifold is orientable, two-sided and
orientable components coincide, and the paired double covers are
orientable. If $O$ and $N$ are the numbers of orientable and nonorientable
components of $y$, the scaled surface has
\[
 O_k=kO+\lfloor k/2\rfloor N,\qquad
 N_k=(k\bmod2)N,\qquad C_k=O_k+N_k.
\]
Multiplying the total component count by $k$ would be wrong when $N>0$.
For example, twice a normal M\"obius band is one annulus, not two
M\"obius bands: it has one orientable component, two boundary circles,
Euler characteristic zero and genus zero.

Every boundary circle has an orientable collar in the surface. The normal
orientation cover consequently restricts to two separate copies of each
boundary circle. Thus the total number of boundary circles scales by $k$,
including for one-sided components. Euler characteristic and normal-disc
count also scale by $k$, both from the covering argument and directly from
the linear coordinate formulas. These statements apply to disconnected
surfaces and mixtures of orientable and nonorientable components; no
component-by-component extraction is needed.

\subsection{The producer and the original source}

After full native validation, the producer computes the gcd $k$ of all
$7t$ coordinates. It stops the gcd scan as soon as the value is one. The
empty vector uses the identity case $k=1$. When $k>1$, it divides each
coordinate exactly. Admissibility is preserved: matching equations are
homogeneous, nonnegativity is unchanged, and no new quadrilateral type
appears. Global edge weights, Euler characteristic, normal-disc count and
boundary-arc counts are linear, so their quotient values are obtained by
exact division from the already validated source. There is no second
manifold-validation pass.

The three orbit queries now concern $y$, $2y$ and $\partial y$. Their
counts give $O=C_2-C$ and $N=2C-C_2$ as before. The lifting formula then
returns $O_k$, $N_k$, $C_k$ and $kb$ for the original input $x$.
Connected-surface genus and crosscaps are computed from these lifted counts
and the original Euler characteristic. The original normal-disc count,
maximum coordinate bit length, boundary-arc count and boundary-cohomology
witness are retained. In particular, even $k$ can make the original
boundary class zero although the quotient has nonzero class; using the
quotient witness in its place would be incorrect.

The public \texttt{reduce\_multiplicity} option defaults to true. Setting
it to false retains direct unreduced queries. Primitive and empty inputs
have exactly the previous result structure. A reduced result additionally
reports \texttt{coordinate\_divisor}; its query statistics describe the
quotient, while every topological claim describes the original surface.
The original shared \texttt{max\_cycles} limit applies to the three actual
queries. Cancellation propagates during validation, gcd scanning, division
and orbit work. Incomplete queries remain inconclusive and emit no topology
proof or component count.

\subsection{Versioned proofs without gcd search in replay}

Version-one normal-topology proofs continue to bind and verify the original
three interval systems. The producer retains that format for primitive and
empty inputs and whenever reduction is disabled. A reduced proof uses
\texttt{normal-surface-topology-v2} and supplies its
\texttt{coordinate\_divisor}. Its source digest still binds the full input
$x$, not merely the quotient. The three contained local orbit proofs bind
the reconstructed quotient systems, and its claims concern $x$.

The verifier checks that the supplied divisor is an integer at least two,
that the source is nonempty, and that every source coordinate is divisible
by it. It then reconstructs the quotient geometry and independently checks
all three local traces. Finally it applies the lifting formula and checks
every original topological claim, together with the original boundary
witness. A different divisor with unchanged query traces is rejected
because the reconstructed interval inputs differ.

Replay does not need to prove that the divisor is the greatest common
divisor. Any exact common divisor at least two gives a sound quotient
certificate. For example, a source of multiplicity twelve may use divisor
six and a complete proof for the twice-primitive quotient. This avoids
rerunning producer gcd discovery and supports independently constructed
proofs. The producer chooses the full gcd to minimize the remaining
coordinate sizes. A Boolean, fractional, zero, unit, negative or
nondividing witness is rejected; arbitrary-size hexadecimal integers remain
supported. An empty source cannot masquerade as a version-two reduction.

The event cap remains shared across the three supplied traces. Source
validation and binary division remain necessary work even for very short
traces; a short certificate is not an assertion that the original input
was cheap to read. Version-one proofs from the actual prior implementation
continue to replay with the new checker.

\subsection{Bit costs and proof-size improvement}

Let $B$ bound the original coordinate bit lengths, $B_0$ those after
division, and $t$ the tetrahedron count. Write $P(t,b)$ for the maintained
three-query orbit cost with $O(t)$ pairings and $O(b+\log(t+1))$-bit
endpoints. Validation still processes the full source. In addition, the
reduction uses $O(t)$ gcd and division operations on integers of at most
$B$ bits, followed by
\[
 P(t,B_0)
\]
rather than $P(t,B)$ orbit work. Lifting the fixed number of summary counts
uses ordinary binary multiplication and parity. With schoolbook arithmetic,
a coarse bound for the extra gcd/division work is polynomial in $t$ and
$B$; fast integer arithmetic can reduce that front-end cost. No unit-cost
assumption is made for the supplied large coordinates.

The reduced proof stores the divisor and a fixed number of original
summary counts, together using $O(B+\log(t+1))$ bits, plus the three
quotient traces. Consequently common multiplicity contributes only to the
outer arithmetic fields, rather than to every endpoint in those traces.
For a fixed primitive surface scaled through arbitrarily large integers,
the orbit-search portion and its local-event structure depend only on the
primitive surface. Source decoding, validation, division, output arithmetic
and the outer proof fields continue to depend on the encoded multiplicity.
When the source gcd is one there is no such reduction; the algorithm retains
its previous worst-case polynomial dependence on the input bit length.
This is a parameter-sensitive improvement of a supplied-vector subroutine,
not a new bound on complete normal-vector search.

\subsection{Validation and paired measurements}

The first mathematical check compared the scaling identities with Regina
on 1,080 multiples of the preceding section's 270 actual normal vectors.
These include mixed orientability, disconnected vectors and triangulations
with multiple boundary vertices. The integrated audit additionally compares
complete old and new topology results and both proof formats against the
actual prior producer and checker at commit \texttt{a3f10a6637bb}.

Thirty-six focused tests pass in 7.619 seconds. They cover even and odd
one-sided scaling, connected annuli produced by doubling, mixed components,
20,001-bit multiplicities, original-versus-quotient boundary parity,
nonmaximal but valid divisors, foreign quotient traces, empty and primitive
inputs, strict witness types, exact shared resource boundaries, cancellation
during quotient construction, and replay with gcd and orbit search disabled.
The historical version-one research driver explicitly disables multiplicity
reduction so its original comparison remains reproducible.

All 932 maintained tests pass with Regina in 202.074 seconds. The integrated
native audit completes 1,350 comparisons in 17.493 seconds: 1,090 inputs
use the reduction and 102 resulting surfaces retain a nonorientable
component. Every old proof replays with the current checker, every new
proof verifies, and every topology agrees with Regina and the pinned
producer. Raw scale factors, input-manifest indices, topology, cycle counts
and proof sizes are in \path{data/normal-multiplicity-audit.json}.

The paired benchmark loads the actual prior producer and verifier in
separate module namespaces, retaining unchanged finite geometry and orbit
kernel dependencies. It compares prior count-only calls, current calls with
reduction disabled, current defaults and an identical current control. It
also separately compares old/new production plus verification and old/new
replay of already supplied proofs. Each fixture has one shuffled warmup and
five measured rounds, with small inputs batched. Every computation includes
fresh source validation. Serialization is outside timing; final topology
and certificates are checked against the reference after the measured
calls. The combined proof arms include verification inside timing.

\begin{center}\small
\begin{tabular}{@{}lrrrrr@{}}
\toprule
Input & old ms & new ms & count ratio & proof ratio & proof bytes old/new \\
\midrule
primitive1 & 0.505 & 0.514 & 0.983 & 0.994 & 2820/2820 \\
primitive8 & 4.948 & 4.976 & 0.995 & 1.017 & 20387/20387 \\
primitive32 & 70.897 & 71.281 & 0.979 & 0.990 & 234409/234409 \\
primitive96 & 1218.134 & 1217.641 & 1.000 & 0.989 & 2018995/2018995 \\
parallel5000 & 9.821 & 5.109 & 1.914 & 1.692 & 434035/27931 \\
parallel20000 & 21.757 & 5.671 & 3.889 & 2.987 & 1671535/50431 \\
one\_sided20000 & 0.738 & 0.352 & 2.102 & 1.779 & 231732/26445 \\
mixed20000 & 3.152 & 0.968 & 3.283 & 2.434 & 620065/34638 \\
\bottomrule
\end{tabular}
\end{center}


Count and certified ratios are medians of paired old/new times. The former
compares the full supplied-vector topology call; the latter includes proof
production and replay on both sides. Neither includes finding a normal
vector or constructing its source manifold from a knot diagram. Exact
representative old/new proofs, all samples, query statistics and loaded
source hashes are retained in
\path{../fast/results/normal_multiplicity_20261008.json}.

The principal run contains 1,560 measured calls and 312 excluded warmup
calls and takes 59.287 seconds. For $2^{20000}$ parallel copies of the
eight-tetrahedron meridian, the complete count-only median falls from
21.757 to 5.671 ms, with median paired ratio 3.889. Production plus replay
falls from 33.796 to 11.313 ms (ratio 2.987), and compact proof transport
falls from 1,671,535 to 50,431 bytes, about a 33-fold reduction.
For $2^{5000}$ copies, count and certified ratios are 1.914 and 1.692,
with proof size 434,035 versus 27,931 bytes.

The odd one-sided multiplicity $2^{20000}+1$ has count and certified ratios
2.102 and 1.779; its proof shrinks from 231,732 to 26,445 bytes. A compatible
mixture of a vertex-linking disc and a M\"obius band at the same scale has
ratios 3.283 and 2.434, with 620,065 versus 34,638 proof bytes. These are
complete supplied-vector computations, rather than timings that exclude
source validation or correctness checks from one side of the comparison.

Primitive controls have paired count ratios from 0.979 to 1.000 and retain
identical proof sizes. The identical current controls range from 0.974 to
1.019 across the principal run. No benefit is claimed for primitive inputs;
the gcd scan adds a small amount of work and quickly stops when it reaches
one. The timing differences on those controls are much smaller than the
large-multiplicity effects.

A separate small-multiplicity follow-up uses four families---one- and
eight-tetrahedron meridians, the one-sided core, and a mixed disc/band
surface---at scales $2$, $3$, $8$, $64$ and $2^{32}$. Seven shuffled
measured rounds of ten calls per arm produce 5,600 timed calls, with 800
excluded warmups. Old versus identical-old control ratios range from
0.969 to 1.021. Count-only old/new ratios range from 0.995 to 1.179.
The twice-meridian at one tetrahedron is essentially unchanged
(0.488/0.492 ms), while the eight-tetrahedron $2^{32}$ case improves from
5.477 to 4.817 ms (paired ratio 1.123). Every one of these twenty cases
finishes with the previous cycle allowance. Raw samples and cap probes are
in \path{../fast/results/normal_multiplicity_small_20261008.json}.

Neither runtime nor proof size is asserted to improve on every possible
input. In the 1,350-case native audit, 58 small certificates grow rather
than shrink; the largest growth is from 12,883 to 13,405 bytes. The divisor
field and differences in local reduction traces can outweigh endpoint
savings at small scale. No audited case uses more orbit cycles after
reduction, but that finite comparison is not a theorem of cycle-count
monotonicity for the scheduler. The unchanged direct path remains available
for callers and future adaptive scheduling experiments.

\paragraph{Complexity target.}
This change isolates a large numeric factor and preserves a proof for the
original object. It removes common multiplicity from the expensive local
orbit stage, while charging full-source arithmetic. It does not supply a
normal vector, certify correspondence with the input knot exterior, or
bound a complete hierarchy search. Those obligations remain necessary for
a general quasi-polynomial unknot recognizer.
